%%=============================================================================
%% Proof-of-Concept
%%=============================================================================

\chapter{\IfLanguageName{dutch}{Proof-of-Concept}{Proof of Concept}}%
\label{ch:poc}

Dit hoofdstuk beschrijft de praktische uitwerking van het onderzoek. De nadruk ligt op de uiteindelijke CTFd-plugin, omdat die rechtstreeks aansluit bij de onderzoeksvraag: challenge-containers dynamisch beheren vanuit CTFd. Voorafgaand daaraan werd echter een beperkte standalone PoC gebouwd om de kernmechanismen sneller te testen. Die verkennende stap wordt kort besproken omdat ze verschillende ontwerpkeuzes voor de plugin heeft onderbouwd.

%%----------------------------------------------------------------------
\section{Architectuuroverzicht}%
\label{sec:poc-architectuur}

De uiteindelijke PoC bestaat uit een CTFd-instantie waarin een eigen plugin een nieuw challenge-type toevoegt. Die plugin spreekt de Docker-daemon aan via de Docker SDK voor Python~\autocite{DockerSDKPython} en beheert de volledige levenscyclus van per-gebruiker of per-team challenge-instanties. Dezelfde plugin kan zowel een lokale Docker-daemon via de Unix-socket aanspreken als een remote Docker API op een andere host. Runtime-metadata, zoals actieve instanties en image-archieven, worden lokaal bijgehouden zodat de plugin ook cleanup en administratieve acties kan uitvoeren.

De opzet bevat vier kernonderdelen:

\begin{itemize}
  \item \textbf{CTFd}: het wedstrijdplatform met authenticatie, challenges en scoreverwerking.
  \item \textbf{CTFd-plugin}: de eigenlijke orkestratielaag die start-, stop- en cleanupacties uitvoert.
  \item \textbf{Docker-daemon}: de runtime waarin challenge-containers en hun netwerken worden aangemaakt; in split-host modus draait die op een afzonderlijke host.
  \item \textbf{Runtime-opslag}: een lichte SQLite-opslag voor actieve instanties, logs en metadata van geuploade image-archieven.
\end{itemize}

Tijdens de implementatie werd eerst vertrokken van een single-host opstelling om de integratie beheersbaar te houden. Nadien is dezelfde architectuur uitgebreid naar een split-host model waarin het Docker-control endpoint en de spelergerichte runtime-host afzonderlijk configureerbaar zijn. Binnen de validatie van deze bachelorproef werd die tweede modus zowel lokaal gesimuleerd als op twee afzonderlijke VM's getest.

\begin{figure}[H]
  \centering
  \resizebox{\textwidth}{!}{%
  \begin{tikzpicture}[
    node distance=1.05cm and 1.25cm,
    font=\small,
    box/.style={draw, rounded corners, align=center, text width=2.75cm, minimum height=1.25cm, fill=black!3},
    plugin/.style={draw, rounded corners, align=center, text width=3.2cm, minimum height=1.7cm, fill=black!6},
    store/.style={draw, rounded corners, align=center, text width=3.2cm, minimum height=1.25cm, fill=white},
    arrow/.style={->, thick},
    note/.style={font=\scriptsize, align=center, fill=white, inner sep=1.5pt}
  ]
    \node[box] (player) {\textbf{Deelnemer of\\ organisator}\\start, stop, solve\\beheeracties};
    \node[box, right=of player] (ctfd) {\textbf{CTFd}\\challenge-flow\\beheerpagina};
    \node[plugin, right=of ctfd] (plugin) {\textbf{Containerized plugin}\\policy- en lifecyclelogica\\plugin-hooks\\views};
    \node[box, right=of plugin] (docker) {\textbf{Docker-daemon}\\containers\\netwerken};
    \node[box, below=of docker] (runtime) {\textbf{Challenge-instantie}\\per gebruiker\\of team};
    \node[store, below=of plugin] (storage) {\textbf{Runtime-opslag}\\SQLite\\image-archieven\\logs en metadata};

    \draw[arrow] (player) -- node[note, above=3pt, pos=0.6] {start / stop / solve} (ctfd);
    \draw[arrow] (ctfd) -- node[note, above=3pt, pos=0.58] {plugin-hooks\\views} (plugin);
    \draw[arrow] (plugin) -- node[note, above=3pt, pos=0.55] {Docker SDK} (docker);
    \draw[arrow] (docker) -- node[note, right=3pt] {run / rm\\network create} (runtime);
    \draw[arrow] (plugin) -- node[note, left=3pt] {metadata\\image-archieven} (storage);
  \end{tikzpicture}
  }
  \caption[Architectuur van de plugin-PoC]{Architectuur van de plugin-PoC. CTFd blijft het centrale platform, terwijl de plugin de lifecycle van challenge-containers en de bijhorende metadata beheert. De Docker-daemon kan lokaal of op een afzonderlijke runtime-host draaien.}
  \label{fig:poc-architectuur}
\end{figure}

%%----------------------------------------------------------------------
\section{Verkennende standalone PoC}%
\label{sec:poc-standalone}

Voor de plugin werd uitgewerkt, werd eerst een kleine standalone orchestrator gebouwd in Flask. Het doel daarvan was niet om een alternatieve eindarchitectuur naar voren te schuiven, maar om drie vragen snel te beantwoorden:

\begin{itemize}
  \item Kan een Python-service containers starten, stoppen en opruimen via de Docker SDK?
  \item Welke basisbeveiligingsmaatregelen zijn praktisch toepasbaar zonder de challenge-flow onnodig complex te maken?
  \item Hoe kunnen lifecycle-regels zoals timeouts, quota en idempotent starten eenvoudig gevalideerd worden?
\end{itemize}

De standalone PoC bood een challenge-registry, een eenvoudige webinterface, API-endpoints voor start en stop, een timeout-reaper en een beperkte logginglaag. Uit die stap volgden een aantal keuzes die rechtstreeks in de plugin zijn meegenomen:

\begin{itemize}
  \item gebruik van de Docker SDK in plaats van shell-calls naar de Docker CLI;
  \item idempotent startgedrag voor dezelfde gebruiker of hetzelfde team;
  \item standaardresourceprofielen per challenge;
  \item cleanup van containers en netwerken bij timeout of expliciete stop;
  \item netwerkisolatie per instantie.
\end{itemize}

De standalone PoC was dus vooral een technische voorstudie die toeliet om de risicovolste runtime-aspecten vroeg te toetsen. Zodra die basis werkte, werd de focus verschoven naar de CTFd-plugin, omdat die nauwer aansluit bij het doel van dit onderzoek: de containerlevenscyclus rechtstreeks koppelen aan acties in het CTF-platform.

\begin{table}[htbp]
  \centering
  \caption[Vergelijking van beide PoC's]{Vergelijking tussen de verkennende standalone PoC en de uiteindelijke CTFd-plugin PoC.}
  \label{tab:poc-vergelijking}
  \begin{tabularx}{\textwidth}{@{}lXX@{}}
    \toprule
    \textbf{Aspect} & \textbf{Standalone PoC} & \textbf{CTFd-plugin PoC} \\
    \midrule
    Hoofddoel & Snelle validatie van lifecycle-mechanismen en beveiligingsmaatregelen. & Integratie van die mechanismen in de echte challenge-flow van CTFd. \\
    Interface & Eigen Flask-dashboard en REST-API. & Challenge-type en beheerpagina binnen CTFd. \\
    Bewijswaarde & Toont dat Docker-orkestratie technisch werkt in isolatie. & Toont dat gebruikersacties in CTFd rechtstreeks runtimegedrag kunnen sturen. \\
    Belangrijkste output & Inzichten over idempotent starten, timeout cleanup en isolatie. & Werkende plugin met uploadflow, solve-cleanup en administratieve controle. \\
    Rol in deze bachelorproef & Verkennende tussenstap om ontwerpkeuzes te onderbouwen. & Hoofdresultaat dat de onderzoeksvraag het dichtst benadert. \\
    \bottomrule
  \end{tabularx}
\end{table}

%%----------------------------------------------------------------------
\section{Implementatie van de CTFd-plugin}%
\label{sec:poc-plugin}

\subsection{Custom challenge-type en configuratie}

De plugin voegt een nieuw challenge-type \texttt{containerized} toe aan CTFd. Een organisator kan zo per challenge runtime-parameters configureren zonder externe scripts of manuele \texttt{docker run}-commando's. De belangrijkste parameters zijn:

\begin{itemize}
  \item Docker-image of geupload image-archief;
  \item een of meerdere containerpoorten;
  \item CPU-limiet;
  \item geheugenlimiet;
  \item timeout;
  \item maximum aantal gelijktijdige instanties.
\end{itemize}

Naast een klassieke image-referentie ondersteunt de plugin ook het uploaden van een \texttt{.tar}, \texttt{.tar.gz} of \texttt{.tgz}-archief uit \texttt{docker save}. Dat archief wordt lokaal opgeslagen, in de Docker-daemon geimporteerd en aan de challenge gekoppeld. Daardoor kan een organisator ook zonder externe registry een image beschikbaar maken voor de PoC. Wanneer de lokaal geimporteerde image later ontbreekt, kan de plugin het gekoppelde archief automatisch opnieuw importeren bij een volgende runtime-start. Die eigenschap bleek bijzonder nuttig in de smoke tests, omdat het daarmee mogelijk werd om de volledige upload- en herimportflow reproduceerbaar te valideren.

\subsection{Start- en stopflow voor deelnemers}

Wanneer een deelnemer een containerized challenge opent, kan die via het challenge-scherm een eigen runtime-instantie starten. De plugin vertaalt die actie naar een Docker-start met de juiste challenge-configuratie. In individuele modus gebeurt dit per gebruiker; in teammodus per team. Daarbij gelden twee belangrijke beleidsregels:

\begin{itemize}
  \item voor dezelfde gebruiker of hetzelfde team wordt een bestaande actieve instantie hergebruikt in plaats van een nieuwe te starten;
  \item zodra het ingestelde maximum aantal instanties is bereikt, weigert de plugin bijkomende starts.
\end{itemize}

Na een succesvolle start krijgt de deelnemer een of meerdere toegangspunten terug met de gepubliceerde hostpoort(en). Die URL's worden opgebouwd op basis van een publiek configureerbare runtime-host, zodat de plugin ook correct blijft werken wanneer CTFd en Docker niet op dezelfde machine draaien. De plugin bewaart tegelijk metadata zoals eigenaar, starttijd, vervaltijd, container-ID, netwerknaam en poortbindingen. Wanneer de deelnemer expliciet stopt, wordt de container verwijderd en wordt ook het bijhorende Docker-netwerk opgeruimd.

\subsection{Timeout cleanup en cleanup bij solve}

De plugin bevat een achtergrondproces dat periodiek controleert welke instanties vervallen zijn. Wanneer een timeout bereikt is, wordt de container geforceerd gestopt en worden de runtime-resources vrijgemaakt. De eindstatus van zo'n instantie wordt als verlopen geregistreerd zodat achteraf zichtbaar blijft waarom ze stopte.

Daarnaast grijpt de plugin ook in op het moment dat een challenge correct opgelost wordt. In dat geval wordt de actieve instantie van de betrokken gebruiker of het betrokken team automatisch stopgezet. Die koppeling is belangrijk voor CTF-contexten: een deelnemer hoeft een oplossing niet eerst manueel te stoppen en de infrastructuur blijft niet onnodig resources verbruiken nadat de flag al werd ingediend.

\begin{figure}[H]
  \centering
  \resizebox{0.95\textwidth}{!}{%
  \begin{tikzpicture}[
    node distance=0.9cm and 0.95cm,
    flowbox/.style={draw, rounded corners, align=center, minimum width=2.65cm, minimum height=0.85cm, fill=black!3},
    arrow/.style={->, thick}
  ]
    \node[flowbox] (open) {Challenge\\ openen};
    \node[flowbox, right=of open] (policy) {Quota en\\ timeout check};
    \node[flowbox, right=of policy] (start) {Container\\ starten};
    \node[flowbox, right=of start] (access) {Runtime-URL\\ tonen};

    \node[flowbox, below=of access] (active) {Actieve\\ instantie};
    \node[flowbox, below left=of active] (stop) {Stop of\\ solve};
    \node[flowbox, below right=of active] (timeout) {Timeout\\ bereikt};
    \node[flowbox, below=1.05cm of active] (cleanup) {Container en\\ netwerk opruimen};

    \draw[arrow] (open) -- (policy);
    \draw[arrow] (policy) -- (start);
    \draw[arrow] (start) -- (access);
    \draw[arrow] (access) -- (active);
    \draw[arrow] (active) -- (stop);
    \draw[arrow] (active) -- (timeout);
    \draw[arrow] (stop) -- (cleanup);
    \draw[arrow] (timeout) -- (cleanup);
  \end{tikzpicture}
  }
  \caption[Start-, stop- en cleanupflow]{Start-, stop- en cleanupflow voor een containerized challenge. Dezelfde cleanupstap wordt gebruikt bij manueel stoppen, een correcte solve en timeout.}
  \label{fig:runtime-flow}
\end{figure}

\subsection{Administratieve runtime-interface}

Voor organisatoren voorziet de plugin een aparte beheerpagina binnen CTFd. Daarop zijn twee soorten informatie zichtbaar:

\begin{itemize}
  \item een overzicht van alle containerized challenges met hun resourceprofiel;
  \item een runtime-overzicht van actieve en historische instanties.
\end{itemize}

Vanuit die interface kunnen beheerders een individuele instantie geforceerd stoppen, alle actieve instanties van een challenge stoppen en de timeout-reaper manueel uitvoeren. Dit verhoogt de beheerbaarheid van de PoC aanzienlijk, omdat de organisator niet rechtstreeks op de Docker-host hoeft in te loggen om runtime-problemen op te lossen.

\clearpage

%%----------------------------------------------------------------------
\section{Beveiliging in de plugin-PoC}%
\label{sec:poc-beveiliging}

De plugin past bij elke containerstart een aantal direct toepasbare basismaatregelen toe, in lijn met de OWASP Docker Security Cheat Sheet~\autocite{OWASP2024} en het CIS Docker Benchmark~\autocite{CISDocker}. Concreet gaat het om de volgende Docker-instellingen:

\begin{itemize}
  \item \textbf{Read-only root filesystem}: het hoofdbestandssysteem van de container wordt alleen-lezen gemount. De challenge kan daardoor niet zomaar systeembestanden, configuratie of applicatiebestanden in de container aanpassen. Dat beperkt de impact van een exploit die probeert persistentie of extra tooling in de container weg te schrijven.
    \item \textbf{\texttt{cap\_drop=['ALL']}}: Linux-capabilities zijn fijnmazige rechten die processen extra kernelmogelijkheden geven, zoals netwerkconfiguratie aanpassen of systeembeheeracties uitvoeren. Door alle capabilities standaard weg te nemen, start de container met zo weinig mogelijk privileges. Alleen wanneer een challenge aantoonbaar extra rechten nodig heeft, zou een organisator daar bewust van kunnen afwijken.
  \item \textbf{\texttt{security\_opt=['no-new-privileges:true']}}: deze instelling verhindert dat processen binnen de container via mechanismen zoals setuid-binaries extra privileges verkrijgen. Zelfs wanneer een programma binnen de container normaal gezien met verhoogde rechten zou starten, mag het proces geen nieuwe privileges krijgen bovenop wat het al had.
  \item \textbf{PID-limiet}: het aantal processen dat binnen de container mag draaien wordt begrensd. Dit voorkomt eenvoudige fork bombs of challenges die per ongeluk zeer veel processen aanmaken en zo de Docker-host belasten.
  \item \textbf{CPU- en geheugenlimieten}: per challenge wordt vastgelegd hoeveel CPU-tijd en RAM een container maximaal mag gebruiken. Daardoor kan één foutieve of kwaadwillige challenge-instantie niet alle resources van de host opeisen en andere deelnemers of CTFd zelf verstoren.
  \item \textbf{Tijdelijke \texttt{tmpfs}-mount voor \texttt{/tmp}}: de map \texttt{/tmp} blijft beschrijfbaar, maar wordt in geheugen geplaatst en verdwijnt wanneer de container stopt. Zo kunnen applicaties die tijdelijke bestanden nodig hebben blijven werken, terwijl permanente writes naar het containerbestandssysteem vermeden worden.
  \item \textbf{Bridge-netwerk per runtime-instantie}: elke challenge-instantie krijgt een eigen Docker-netwerk. Containers van verschillende gebruikers of teams delen daardoor geen standaard netwerksegment en kunnen elkaar niet zomaar via containernaam of intern IP-adres ontdekken.
\end{itemize}

Die combinatie beperkt het risico op privilege-escalatie, laterale beweging en resource-uitputting. De keuze voor per-instantie netwerken zorgt ervoor dat deelnemers niet zomaar andere challenge-instanties kunnen ontdekken of aanspreken. Voor split-host deployments ondersteunt de plugin daarnaast een expliciet configureerbaar gepubliceerd poortbereik, zodat firewallregels op de runtime-host begrensd kunnen blijven. Belangrijk is wel dat deze PoC een basisbeveiliging demonstreert en geen volledig production-ready hardeningprofiel oplevert. Zo is ondersteuning voor aangepaste seccomp-profielen nog geen afgewerkt onderdeel van de implementatie en werd remote Docker-communicatie in de validatie enkel op wegwerpomgevingen zonder TLS getest.

\clearpage

%%----------------------------------------------------------------------
\section{Deployment- en schaalbaarheidsoverwegingen}%
\label{sec:poc-schaalbaarheid}

De huidige plugin-PoC ondersteunt twee deploymentmodi. In same-host modus draait CTFd samen met de plugin op dezelfde host als de Docker-daemon en wordt de lokale Unix-socket gebruikt. In split-host modus blijft CTFd het control plane, maar wordt de Docker-daemon op een andere host aangesproken via een afzonderlijk control endpoint. De spelergerichte runtime-host kan daarbij verschillen van het Docker API endpoint. Binnen die scope toont de PoC aan dat een plugin-gebaseerde aanpak haalbaar is voor:

\begin{itemize}
  \item dynamisch starten en stoppen van challenge-containers;
  \item afdwingen van limieten per challenge;
  \item automatisch cleanupgedrag bij timeout en solve;
  \item administratieve controle vanuit CTFd;
  \item scheiding tussen control plane en runtime-host in een split deployment.
\end{itemize}

De split-host ondersteuning werd in deze bachelorproef ook effectief gevalideerd. Eerst gebeurde dat lokaal via een Compose-opstelling met een aparte \texttt{docker:dind}-service en een smoke-runner binnen hetzelfde netwerk. Daarna volgde een reële validatie op twee aparte Lima-VM's: één VM voor CTFd en de plugin, en één VM voor de challenge-containers. Daarbij werden de gepubliceerde challenge-poorten begrensd tot een vast bereik, zodat de runtime-host eenvoudiger afgeschermd kon worden.

Voor verdere schaalvergroting blijven wel nog een aantal vragen over:

\begin{itemize}
  \item hoe remote Docker-communicatie in productie TLS-beveiligd en operationeel beheersbaar wordt ingericht;
  \item hoe het systeem reageert wanneer de runtime-host tijdelijk onbereikbaar wordt tijdens start, stop of cleanup;
  \item hoe dynamische routing best wordt gecombineerd met een reverse proxy in plaats van louter hostpoortmapping;
  \item hoe meerdere runtime-hosts evenwichtig kunnen worden aangestuurd vanuit dezelfde plugin.
\end{itemize}
